% =============================================================================
% Paper Transport Conjecture -- Explicit formal statement + four proof roadmaps
% Target : arXiv math-ph (followed by submission to Comm. Math. Phys. or similar)
% Author : Kevin Remondiere (Oloron-Sainte-Marie, France)
% Status : 4-6 pp, TIER 3 OPEN problem statement, honest no-overclaim
% =============================================================================
\documentclass[11pt,a4paper,reqno]{amsart}

\usepackage[utf8]{inputenc}
\usepackage[T1]{fontenc}
\usepackage{amsmath,amssymb,amsthm,mathrsfs}
\usepackage{geometry}
\usepackage{hyperref}
\usepackage{enumitem}
\usepackage{booktabs}
\usepackage{xcolor}
\usepackage{microtype}

\geometry{margin=2.5cm}

\hypersetup{
  colorlinks=true,
  linkcolor=blue!60!black,
  citecolor=blue!60!black,
  urlcolor=blue!60!black
}

\newtheorem{theorem}{Theorem}[section]
\newtheorem{proposition}[theorem]{Proposition}
\newtheorem{lemma}[theorem]{Lemma}
\newtheorem{corollary}[theorem]{Corollary}
\newtheorem{conjecture}[theorem]{Conjecture}
\theoremstyle{definition}
\newtheorem{definition}[theorem]{Definition}
\newtheorem{remark}[theorem]{Remark}
\newtheorem{strategy}{Strategy}

\DeclareMathOperator{\PSL}{PSL}
\DeclareMathOperator{\SL}{SL}
\DeclareMathOperator{\SU}{SU}
\DeclareMathOperator{\Cl}{Cl}
\DeclareMathOperator{\Vol}{Vol}
\newcommand{\OK}{\mathcal{O}_K}
\newcommand{\YK}{Y_K}
\newcommand{\KASP}{K_{\mathrm{ASP}}}
\newcommand{\HH}{\mathbb{H}}
\newcommand{\ZZ}{\mathbb{Z}}
\newcommand{\QQ}{\mathbb{Q}}
\newcommand{\RR}{\mathbb{R}}
\newcommand{\CC}{\mathbb{C}}
\newcommand{\NN}{\mathbb{N}}
\newcommand{\TT}{\mathbb{T}}

\title[Transport Conjecture]
  {The Transport Conjecture between lattice SU($N$) glueballs and Bianchi
   orbifold spectral observables: formal statement and four candidate proof
   strategies}

\author[K.~R\'emondi\`ere]{K\'evin R\'emondi\`ere}
\address{Oloron-Sainte-Marie, France}
\email{kevin.remondiere@gmail.com}
\thanks{ORCID 0009-0008-2443-7166. This work is part of the
\emph{crossed-cosmos} project (concept DOI 10.5281/zenodo.19686398).}

\subjclass[2020]{Primary 81T13, 81T25; Secondary 11F72, 58J50}

\keywords{Yang--Mills, lattice gauge theory, Bianchi 3-orbifold, Wightman,
constructive QFT, Selberg pretrace, Transport Conjecture}

\date{\today}

\begin{document}

\begin{abstract}
We formulate explicitly an open conjectural identification, named the
\emph{Transport Conjecture}, between (i) the lattice SU($N$) scalar glueball
mass $m_{0^{++}}(\SU(N))$, and (ii) a spectral observable
$m^{(\chi_*)}_{\mathrm{spectral}}$ on a Bianchi 3-orbifold $\YK$ with
$K = \KASP(N)$ a privileged imaginary quadratic anchor. The conjecture is
needed as an explicit substep in a surrogate mass-gap formula proposed in the
companion paper \cite{Remondiere2026RouteB}; the present preprint isolates
the conjecture as an \emph{open problem} (TIER~3 in our internal
classification), gives its explicit formal statement, and lists four
candidate proof strategies (CCHS 4D extension; Cao--Sheffield random surfaces;
Wightman positivity W0; Cluster decomposition W5). We also state five
falsifiable per-anchor predictions, and conclude with an honest assessment
of the open-problem status.
\end{abstract}

\maketitle

% =============================================================================
\section{The Transport Conjecture}\label{sec:conjecture}
% =============================================================================

\subsection{Setup}\label{ss:setup}

Let $N\ge 2$ be an integer and let $\TT^4 = (\RR/L\ZZ)^4$ be the 4-torus of edge $L$. Let $\mathcal{A}_{\SU(N)}(\TT^4)$ denote the lattice SU($N$) Yang--Mills theory on $\TT^4$ with Wilson action (or any of its standard improvements), and let $\mathcal{O}^{(0^{++})} : \mathcal{A}_{\SU(N)}(\TT^4) \to \RR$ denote the lowest scalar glueball observable. In the continuum limit (lattice spacing $a \to 0$ at fixed physical volume), one defines
\begin{equation}\label{eq:mlat}
   m_{0^{++}}^{\mathrm{lat}}(\SU(N))
   \;:=\; \lim_{a\to 0}\, \sqrt{a}\cdot\bigl\langle \mathcal{O}^{(0^{++})}\,\hat{\mathcal{O}}^{(0^{++})}\bigr\rangle^{1/2}_{a},
\end{equation}
where the $\hat{}$ denotes a temporal-correlator extraction and $\langle\cdot\rangle_a$ the lattice ensemble average at spacing $a$.

On the arithmetic side, let $K=\QQ(\sqrt{D})$ be imaginary quadratic with fundamental discriminant $D<0$, $\OK$ its ring of integers, $\YK := \PSL_2(\OK)\backslash\HH^3$ the associated Bianchi 3-orbifold \cite[\S 1.5]{ElstrodtGrunewaldMennicke1998}, and $\Delta$ the positive hyperbolic Laplacian on $\YK$. By the per-character spectral decomposition of the companion paper \cite{Remondiere2026Selberg},
\[
   L^2(\YK) = \bigoplus_{\chi \in \widehat{g}(K)} L^2(\YK)_\chi,
\]
and $\Delta_\chi := \Delta\big|_{L^2(\YK)_\chi}$ has its own discrete cuspidal + continuous Eisenstein spectrum. The Center--Rank theorem CR of \cite{Remondiere2026CenterRank} selects a privileged genus character $\chi_*$ for the anchor $\KASP(N)$ (the ``central'' character compatible with the 2-part of $Z(\SU(N))$). Define
\begin{equation}\label{eq:mspectral}
   m^{(\chi_*)}_{\mathrm{spectral}}(K, N)
   \;:=\; \sqrt{\lambda_1(\Delta_{\chi_*})}\,,
\end{equation}
the square root of the lowest cuspidal eigenvalue of $\Delta_{\chi_*}$ on $\YK$.

\subsection{Conjecture}\label{ss:conj}

\begin{conjecture}[Transport Conjecture]\label{conj:transport-formal}
Let the universal constants $\sqrt{2\pi e}$, $\sqrt{2/3}$ and the Dijkgraaf--Witten factor $F(N) = (9/10)(N^2+1)/N^2$ be as in \cite{Remondiere2026RouteB}. Let
\[
   \mathcal{C}(N) \;:=\; \sqrt{2\pi e}\,\cdot\,\sqrt{\tfrac{2}{3}}\,\cdot\,F(N)\,.
\]
There exists a Bianchi orbifold $\YK$ with $K = \KASP(N)$ (in the sense of the Anchor Selection Principle of \cite{Remondiere2026CenterRank}) such that, in the lattice continuum limit,
\begin{equation}\label{eq:transport}
   \frac{m_{0^{++}}^{\mathrm{lat}}(\SU(N))}{\sqrt{\sigma_N}}
   \;=\; \mathcal{C}(N)\,,
\end{equation}
where $\sqrt{\sigma_N}$ is the SU($N$) string tension (extracted from the lattice via the standard Wilson loop area-law, or by any equivalent prescription).

Equivalently, there exists a \emph{transport map} $\Phi : \mathcal{O}^{(0^{++})}|_{\TT^4} \mapsto m^{(\chi_*)}_{\mathrm{spectral}}(\KASP(N), N)$ such that the right-hand side of \eqref{eq:transport} factorises through the spectral observable on $\YK$.
\end{conjecture}

\subsection{Numerical status (2026)}\label{ss:numeric}

The Athenodorou--Teper 2021 lattice continuum data \cite[Tab.~34]{AthenodorouTeper2021} for SU($N$), $N\in\{2,3,4,5,6,\infty\}$, gives an RMS deviation of $0.7\%$ from the conjectural value $\mathcal{C}(N)$ defined in \eqref{eq:transport}; see Table~1 of \cite{Remondiere2026RouteB}. This numerical match is consistent with, but does not prove, the conjecture.

% =============================================================================
\section{Four candidate proof strategies}\label{sec:strategies}
% =============================================================================

We list four candidate strategies for promoting Conjecture~\ref{conj:transport-formal} from TIER 3 (explicit open conjecture, with $0.7\%$ RMS empirical match) to TIER 1 (rigorous proof). Each strategy is treated as a programmatic outline, not a road-mapped proof.

\begin{strategy}[CCHS 4D extension]\label{strat:CCHS}
The recent work of Chatterjee, Cao, Hairer, and Shen (\cite{Chatterjee2020} and follow-ups; the 4D extension is, to the author's knowledge, still in progress as of mid-2026) constructs lattice SU($N$) Wightman observables in $2$D and $3$D via regularity-structures techniques. A 4D extension, when it exists, would provide a rigorous existence theorem for $m_{0^{++}}^{\mathrm{lat}}(\SU(N))$ as a Wightman 2-point function correlation. The proof of Conjecture~\ref{conj:transport-formal} via this route would then require: (a) the 4D CCHS construction; (b) identification of the Wightman 2-point function with the spectral observable on a Bianchi orbifold via a Bianchi-cycle gauge fixing.

\textbf{Status (2026):} CCHS-3D published; CCHS-4D in progress (no precise prediction by the author for completion date). Bianchi-cycle identification: open. Pessimistic 15-year estimate: $10$-$20\%$.
\end{strategy}

\begin{strategy}[Cao--Sheffield random surfaces]\label{strat:CaoSheffield}
A recent programme by Cao and Sheffield (and collaborators) constructs random surface (Liouville) models that, in a certain limit, factor through gauge-theory observables on $\TT^4$. If this construction extends to Bianchi 3-orbifolds (a manifold without flat 4D structure), it could provide a Bianchi-to-lattice transport: the Wilson loops on $\TT^4$ would arise as random-surface measure averages, and the Bianchi-orbifold spectral observables would be obtained as a homotopy quotient of the Liouville model.

\textbf{Status (2026):} 2D Cao--Sheffield programme partly published; extension to $4$D and to non-compact Bianchi 3-orbifolds is open.
\end{strategy}

\begin{strategy}[Wightman positivity (W0)]\label{strat:Wightman}
Direct Wightman positivity \cite{StreaterWightman1964}: prove that the lattice 2-point function in \eqref{eq:mlat} is the spectral measure of an appropriate self-adjoint operator on a Hilbert space, and that this operator coincides with $\Delta_{\chi_*}$ on $L^2(\YK)_{\chi_*}$ after identification of the gauge-invariant Hilbert space with the Bianchi spectral Hilbert space. This is, in essence, the constructive QFT roadmap for Yang--Mills mass gap, restricted to the surrogate observable.

\textbf{Status (2026):} The full Wightman programme for 4D Yang--Mills is an open Clay Millennium problem \cite{JaffeWitten2000}. A restricted version --- proving Wightman positivity \emph{only} for the surrogate observable, not for the full gauge-invariant ensemble --- might be more tractable; this is the author's own preferred long-term route.
\end{strategy}

\begin{strategy}[Cluster decomposition (W5)]\label{strat:Cluster}
Cluster decomposition in the spirit of \cite{Wightman1964}, \cite{HaagKastler1964}: prove that the lattice 2-point function clusters at large separations, hence its spectral support has a gap. If the spectral gap of the lattice theory coincides with the genus character spectral gap of $\Delta_{\chi_*}$ on $\YK$, the Transport Conjecture would follow at the level of spectral support (though not yet at the level of specific eigenvalue).

\textbf{Status (2026):} Cluster decomposition for lattice SU($N$) is folklore-rigorous; the identification with Bianchi spectral support is open.
\end{strategy}

% =============================================================================
\section{Five falsifiable predictions}\label{sec:falsif}
% =============================================================================

The conjecture yields five testable predictions (cf.\ \cite{Remondiere2026RouteB}, \S 7); we restate them here, isolated, for ease of reference. Let $\mathcal{C}(N) := \sqrt{2\pi e}\cdot\sqrt{2/3}\cdot F(N)$.

\paragraph{F-T-1.} $m_{0^{++}}^{\mathrm{lat}}(\SU(2))/\sqrt{\sigma_2} = \mathcal{C}(2) = 3.7962$. \emph{Status (2026):} AT2021 reports $3.781 \pm 0.023$, $\Delta = -0.40\%$, PASS.

\paragraph{F-T-2.} $m_{0^{++}}^{\mathrm{lat}}(\SU(3))/\sqrt{\sigma_3} = \mathcal{C}(3) = 3.3744$. \emph{Status (2026):} AT2021 reports $3.405 \pm 0.021$, $\Delta = +0.91\%$, PASS.

\paragraph{F-T-3.} $m_{0^{++}}^{\mathrm{lat}}(\SU(4))/\sqrt{\sigma_4} = \mathcal{C}(4) = 3.2273$. \emph{Status (2026):} AT2021 reports $3.271 \pm 0.027$, $\Delta = +1.34\%$, PASS.

\paragraph{F-T-4.} $m_{0^{++}}^{\mathrm{lat}}(\SU(5))/\sqrt{\sigma_5} = \mathcal{C}(5) = 3.1584$. \emph{Status (2026):} AT2021 reports $3.156 \pm 0.031$, $\Delta = -0.06\%$, PASS.

\paragraph{F-T-5.} $m_{0^{++}}^{\mathrm{lat}}(\SU(6))/\sqrt{\sigma_6} = \mathcal{C}(6) = 3.1213$. \emph{Status (2026):} AT2021 reports $3.102 \pm 0.032$, $\Delta = -0.61\%$, PASS.

\smallskip

\noindent
The RMS deviation across F-T-1--5 is $0.7\%$, comfortably inside the AT2021 continuum error bars $\sim 3\%$. The empirical evidence is suggestive but not conclusive: no datapoint deviates by more than the $1\sigma$ AT2021 error bar from $\mathcal{C}(N)$. Falsification of any one of F-T-1--5 by future high-precision lattice data would falsify the conjecture.

% =============================================================================
\section{Honest open-problem status}\label{sec:status}
% =============================================================================

The Transport Conjecture is an \emph{open problem}. Its current status is summarised as:
\begin{itemize}
   \item \textbf{Empirically:} consistent with all known lattice data (AT2021, $0.7\%$ RMS over 6 datapoints).
   \item \textbf{Logically:} an explicit identification --- not derivable from the heat-kernel + DW factors alone. Substitutes a non-trivial assumption into the surrogate mass-gap formula of \cite{Remondiere2026RouteB}.
   \item \textbf{Strategically:} all four candidate proof strategies in \S\ref{sec:strategies} are themselves open programmes (CCHS-4D in progress; Cao--Sheffield 2D-only; Wightman positivity is the Clay Millennium problem; Cluster decomposition is folklore but with an open Bianchi identification).
\end{itemize}

We invite the community to attempt the conjecture via any of these (or new) strategies.

\subsection{What this paper does \emph{not} claim}\label{ss:not-claim}

\begin{itemize}
   \item No proof. The conjecture is stated as an open problem.
   \item No claim of solution to the Yang--Mills Clay millennium problem. The conjecture is concerned with a specific surrogate observable.
   \item No claim that the RMS $0.7\%$ empirical match is a proof. Numerical agreement, however suggestive, is not equivalent to a structural argument.
\end{itemize}

% =============================================================================
\section{Acknowledgments}\label{sec:acks}
% =============================================================================

The author thanks the open mathematical literature for the foundational treatments used (Elstrodt--Grunewald--Mennicke, Sarnak, Bunke--Olbrich, Athenodorou--Teper). The work is the result of a structured arithmetic-gauge correspondence programme; full credit for the constituent inputs is given in the companion papers \cite{Remondiere2026Selberg, Remondiere2026CenterRank, Remondiere2026CR, Remondiere2026RouteB}.

% =============================================================================
\bibliographystyle{alpha}
\bibliography{refs}
% =============================================================================

\end{document}
